\documentclass[aspectratio=43]{beamer}

\usetheme{Madrid}
\usecolortheme{default}
\usefonttheme{professionalfonts}
\setbeamertemplate{navigation symbols}{}
\setbeamertemplate{blocks}[rounded][shadow=false]
\setbeamertemplate{items}[circle]

\usepackage{slashed}
\usepackage{datetime}
\usepackage{graphicx}
\usepackage{hyperref}
\usepackage{cancel}
\usepackage[absolute,overlay]{textpos}
\usepackage{tikz}
\usepackage{booktabs}
\usepackage{array}
\usepackage{amsmath,amssymb}
\usepackage{xcolor}

\graphicspath{{plots/}{./}}

\def\Put(#1,#2)#3{\leavevmode\makebox(0,0){\put(#1,#2){#3}}}

\def\mydate{\leavevmode\hbox{\the\year-\twodigits\month-\twodigits\day}}
\def\twodigits#1{\ifnum#1<10 0\fi\the#1}

\newcommand\Wider[2][3em]{%
\makebox[\linewidth][c]{%
  \begin{minipage}{\dimexpr\textwidth+#1\relax}
  \raggedright
  \centering#2
  \end{minipage}%
  }%
}

% ---------- Palette ----------
\definecolor{bg}{HTML}{F7F8FC}
\definecolor{ink}{HTML}{0F172A}
\definecolor{muted}{HTML}{475569}
\definecolor{navy}{HTML}{0B1020}
\definecolor{navy2}{HTML}{11182E}
\definecolor{violet}{HTML}{7C3AED}
\definecolor{violet2}{HTML}{A78BFA}
\definecolor{electric}{HTML}{22D3EE}
\definecolor{rose}{HTML}{F43F5E}
\definecolor{line}{HTML}{D6DAE8}
\definecolor{panel}{HTML}{EEF2FF}

% ---------- Global styling ----------
\setbeamercolor{background canvas}{bg=bg}
\setbeamercolor{normal text}{fg=ink,bg=bg}
\setbeamercolor{frametitle}{fg=white,bg=navy}
\setbeamercolor{title}{fg=white}
\setbeamercolor{subtitle}{fg=violet2}
\setbeamercolor{author}{fg=white}
\setbeamercolor{institute}{fg=violet2}
\setbeamercolor{date}{fg=violet2}

\setbeamercolor{palette primary}{fg=white,bg=navy}
\setbeamercolor{palette secondary}{fg=white,bg=navy2}
\setbeamercolor{palette tertiary}{fg=white,bg=violet}
\setbeamercolor{palette quaternary}{fg=white,bg=electric}

\setbeamercolor{section in toc}{fg=violet}
\setbeamercolor{subsection in toc}{fg=muted}
\setbeamercolor{itemize item}{fg=violet}
\setbeamercolor{itemize subitem}{fg=electric}
\setbeamercolor{enumerate item}{fg=violet}
\setbeamercolor{block title}{fg=white,bg=violet}
\setbeamercolor{block body}{fg=ink,bg=panel}
\setbeamercolor{alerted text}{fg=rose}
\setbeamercolor{caption name}{fg=violet}

\setbeamerfont{title}{series=\bfseries,size=\LARGE}
\setbeamerfont{subtitle}{size=\small}
\setbeamerfont{frametitle}{series=\bfseries,size=\large}
\setbeamerfont{block title}{series=\bfseries}
\setbeamerfont{author}{series=\bfseries}

\hypersetup{
  colorlinks=true,
  linkcolor=violet,
  urlcolor=electric,
  citecolor=rose
}

\setbeamertemplate{frametitle}{
  \nointerlineskip
  \begin{beamercolorbox}[wd=\paperwidth,ht=0.9cm,dp=0.24cm,leftskip=0.42cm,rightskip=0.28cm]{frametitle}
    \usebeamerfont{frametitle}\insertframetitle
  \end{beamercolorbox}
  {\color{electric}\rule{\paperwidth}{0.8pt}}
}

\setbeamertemplate{footline}{
  \leavevmode
  \hbox{
  \begin{beamercolorbox}[wd=.72\paperwidth,ht=0.33cm,dp=0.14cm,leftskip=0.35cm]{author in head/foot}
    {\color{muted}\scriptsize\insertshorttitle}
  \end{beamercolorbox}
  \begin{beamercolorbox}[wd=.28\paperwidth,ht=0.33cm,dp=0.14cm,rightskip=0.35cm]{date in head/foot}
    \hfill{\color{muted}\scriptsize\insertframenumber/\inserttotalframenumber}
  \end{beamercolorbox}}
}

\setbeamertemplate{title page}{
  \begin{tikzpicture}[remember picture,overlay]
    \fill[navy] (current page.south west) rectangle (current page.north east);
    \shade[left color=violet,right color=electric,opacity=0.18]
      ([xshift=-0.8cm,yshift=1.2cm]current page.south east) circle (3.2cm);
    \shade[left color=electric,right color=violet,opacity=0.12]
      ([xshift=0.8cm,yshift=-0.6cm]current page.north west) circle (2.9cm);
    \draw[electric,opacity=0.3,line width=0.7pt]
      ([xshift=0.6cm,yshift=1.0cm]current page.south west) --
      ([xshift=8.3cm,yshift=1.0cm]current page.south west);
  \end{tikzpicture}
  \vspace*{1.35cm}
  \begin{flushleft}
    {\usebeamerfont{title}\usebeamercolor[fg]{title}\inserttitle\par}
    \vspace{0.25cm}
    {\usebeamerfont{subtitle}\usebeamercolor[fg]{subtitle}Computational visualization, prime structure, and experimental mathematics\par}
    \vspace{1.0cm}
    {\usebeamerfont{author}\usebeamercolor[fg]{author}\insertauthor\par}
    \vspace{0.12cm}
    {\usebeamerfont{institute}\usebeamercolor[fg]{institute}\insertinstitute\par}
    \vspace{0.35cm}
    {\usebeamerfont{date}\usebeamercolor[fg]{date}\insertdate\par}
  \end{flushleft}
  \vfill
}

\def\meetingname{PredictiveInsightsAI.com}

\title[\meetingname]{The Riemann Zeta Function:\\Zeros, the Critical Line, and Computational Exploration}
\author[Samuel Castillo]{Samuel Castillo}
\institute[Predictive Insights AI]{www.PredictiveInsightsAI.com}
\date[\mydate]{\meetingname\\\today}

\begin{document}

\begin{frame}[plain]
\titlepage
\end{frame}

\begin{frame}{Roadmap}
\tableofcontents
\end{frame}

\section{Introduction}

\begin{frame}{Introduction}
\begin{block}{From Euler to Riemann}
Euler's resolution of the Basel problem revealed that the zeta function already
encodes remarkable arithmetic structure at special values such as
\[
\zeta(2)=\frac{\pi^2}{6}.
\]
Riemann's later insight extended this perspective dramatically: after analytic
continuation, the non-trivial zeros of the zeta function govern the fine-scale
distribution of prime numbers.
\end{block}
\end{frame}

\begin{frame}{The Riemann Zeta Function}
\begin{itemize}
\item The Riemann zeta function occupies a central place in analytic number theory.
\item For $\Re(s) > 1$, it is defined by the convergent series
\[
\zeta(s)=\sum_{n=1}^{\infty}\frac{1}{n^s}.
\]
\item Euler showed that this function also admits the product expansion
\[
\zeta(s)=\prod_{p}\left(1-p^{-s}\right)^{-1},
\]
which connects analysis directly to the prime numbers.
\end{itemize}
\end{frame}

\begin{frame}{Central Thesis}
\begin{block}{Big idea}
The zeros of $\zeta(s)$ are not merely analytic curiosities. They are structural
features that govern the oscillatory behavior of prime counting.
\end{block}

\medskip
This viewpoint culminates in the Riemann Hypothesis, which asserts that every
non-trivial zero lies on the critical line
\[
\Re(s)=\tfrac{1}{2}.
\]
\end{frame}

\begin{frame}{Prime Number Distribution}
This broader perspective also situates related phenomena in prime number theory.
\begin{itemize}
\item \textbf{Mersenne primes:} numbers of the form $2^p-1$, where $p$ must be prime for $2^p-1$ to have any chance of being prime.
\item \textbf{Twin primes:} pairs of primes differing by $2$, such as $(29,31)$.
\item The Twin Prime Conjecture asserts that infinitely many such pairs exist.
\end{itemize}
\end{frame}

\begin{frame}{Where the Zeros Live}
The non-trivial zeros of the Riemann zeta function lie in the critical strip
\[
0<\Re(s)<1.
\]
Extensive numerical evidence suggests that these zeros lie on the critical line
\[
\Re(s)=\frac{1}{2},
\]
as asserted by the Riemann Hypothesis.

\medskip
Visualization techniques provide useful geometric intuition for the analytic
structure of $\zeta(s)$ in the complex plane.
\end{frame}

\begin{frame}{Visualization Modes}
\begin{itemize}
    \item \textbf{Phase Plots:} hue tracks $\arg(\zeta(s))$ while brightness tracks $|\zeta(s)|$; zeros appear as phase singularities.
    \item \textbf{Contour Plots:} the curves $\Re(\zeta(s))=0$ and $\Im(\zeta(s))=0$ intersect exactly at zeros.
    \item \textbf{Three-Dimensional Surfaces:} the graph of $|\zeta(s)|$ creates a landscape whose minima align with non-trivial zeros.
\end{itemize}
\end{frame}

\begin{frame}{Phase Graph}
\centering
\includegraphics[width=0.79\textwidth]{zeta_phase.png}

\vspace{0.12cm}
{\scriptsize \color{muted}Hue represents $\arg(\zeta(s))$ and brightness represents $|\zeta(s)|$. The dashed line marks $\Re(s)=\tfrac{1}{2}$.}
\end{frame}

\begin{frame}{Surface Plot}
\centering
\includegraphics[width=0.84\textwidth]{z3d.png}

\vspace{0.12cm}
{\scriptsize \color{muted}The minima of $\log(1+|\zeta(s)|)$ correspond to points where $\zeta(s)$ approaches zero.}
\end{frame}

\begin{frame}{Takeaway}
These visual representations reinforce the interpretation of the zeros as
structural features governing the behavior of $\zeta(s)$, and by extension,
the distribution of prime numbers through the explicit formula.
\end{frame}

\section{Computation}

\begin{frame}{Computational Consequences}
\begin{itemize}
    \item \textbf{Computational Evidence:} extensive computations have verified enormous numbers of non-trivial zeros on the critical line.
    \item \textbf{Prime Counting:} the zeros appear explicitly in formulas such as
    \[
    \pi(x)=\mathrm{Li}(x)-\sum_{\rho}\mathrm{Li}(x^\rho)+\text{lower-order terms}.
    \]
    \item \textbf{Error Bounds:} RH implies the near-optimal error term
    \[
    \pi(x)=\mathrm{Li}(x)+O\!\left(x^{1/2}\log x\right).
    \]
\end{itemize}
\end{frame}

\begin{frame}{The Big Formula}
A central bridge between the zeros of $\zeta(s)$ and the distribution of primes
is the explicit formula
\[
\psi(x)=x-\sum_{\rho}\frac{x^\rho}{\rho}-\log(2\pi)-\frac{1}{2}\log(1-x^{-2}),
\]
where $\psi(x)$ is the Chebyshev function and the sum runs over non-trivial
zeros $\rho$ of $\zeta(s)$.

\medskip
This formula shows that the zeros contribute oscillatory corrections to the main term $x$.
\end{frame}

\begin{frame}{Prime Analytics}
\begin{itemize}
    \item \textbf{Numerical contour analysis:} visualize level sets of $\Re(\zeta(s))$ and $\Im(\zeta(s))$, approximate zeros, and study local spacing.
    \item \textbf{Automated theorem search:} machine-assisted systems can explore large proof spaces and intermediate lemmas.
    \item \textbf{Large-scale verification:} distributed computation continues to test zero locations to increasing heights.
\end{itemize}
\end{frame}

\begin{frame}{Classical vs Modern Ciphers}
The contrast between classical and modern cryptography highlights why prime
numbers matter outside pure mathematics.

\begin{block}{Classical example}
A Caesar cipher shifts each letter by a fixed number of positions. A slightly more
complex variant uses a keyword substitution alphabet such as
\[
\text{KEYWORDABCFGHIJLMNPQSTUVXZ}.
\]
\end{block}

These systems are weak because they leak repeated letters, frequency patterns, and key-space structure.
\end{frame}

\begin{frame}{Cipher Comparison}
\small
\centering
\renewcommand{\arraystretch}{1.2}
\begin{tabular}{>{\raggedright\arraybackslash}p{2.25cm}|>{\raggedright\arraybackslash}p{2.45cm}|>{\raggedright\arraybackslash}p{2.5cm}|>{\raggedright\arraybackslash}p{2.35cm}}
\toprule
\textbf{System} & \textbf{Key Space Size} & \textbf{Attack Method} & \textbf{Time to Break} \\
\midrule
Caesar Cipher & 25 keys & Brute force & Instant ($<1$ ms) \\
Keyword Substitution & Large theoretical space & Frequency analysis & Seconds--minutes \\
DES (56-bit) & $2^{56}$ & Exhaustive search & Hours--days \\
RSA-512 & Large composite modulus & Integer factorization & Hours--months \\
RSA-2048 & Very large composite modulus & Integer factorization & Infeasible classically \\
\bottomrule
\end{tabular}
\end{frame}

\begin{frame}{Contour Plot}
\centering
\includegraphics[width=0.76\textwidth]{rzt.png}

\vspace{0.12cm}
{\scriptsize \color{muted}Red curves mark $\Re(\zeta(s))=0$ and green curves mark $\Im(\zeta(s))=0$. Their intersections align near $\Re(s)=\tfrac{1}{2}$.}
\end{frame}

\section{Conclusion}

\begin{frame}{Conclusion}
\begin{itemize}
\item The Riemann Hypothesis links complex analysis with the fine structure of prime distribution.
\item Euler opened one door into the zeta function through special values.
\item Riemann opened another through the zeros, which govern oscillatory structure in the primes.
\end{itemize}
\end{frame}

\begin{frame}{Why the Visuals Matter}
The contour plot makes the geometry concrete:
\begin{itemize}
\item intersections of $\Re(\zeta(s))=0$ and $\Im(\zeta(s))=0$ reveal non-trivial zeros;
\item the explicit formula ties those zeros directly to prime irregularity;
\item computation and visualization create a practical way to explore deep analytic structure.
\end{itemize}
\end{frame}

\begin{frame}{Future Work}
Future work will extend the computational side of this project in four main directions:
\begin{itemize}
    \item high-precision numerical approximation of $\zeta(s)$ along the critical line;
    \item richer visualization, including domain coloring and surface plots;
    \item experiments with truncated explicit formulas for $\psi(x)$ and $\pi(x)$;
    \item exploratory AI-based analysis of zero spacing, clustering, and prime-related data.
\end{itemize}
\end{frame}

\end{document}
